\begin{center}
[5.]
\end{center}
\begin{center}
\textsc{Gauss} an \textsc{Gerling}.\quad Göttingen, 5.~Oktober 1821.
\end{center}

\par\noindent\hspace*{\parindent}Meinen herzlichen Glückwunsch zu der Aufgabe, die Sie übernommen haben.
Es freut mich, dass sie in Ihre Hände kommt, da Sie es sich zur Pflicht machen werden, das ganze Dreiecksnetz mit aller möglichen Sorgfalt auszuführen.
Ich halte das für überaus wünschenswert und bedauere es immer, wenn man schon bei den Dreiecken erster Ordnung knauserig darauf bedacht ist, nur so viel Genauigkeit aufzuwenden, wie für den allerletzten Zweck unumgänglich nötig ist.
Die \textls[120]{genaueste} Kenntnis der relativen Lage der interessantesten Punkte eines Landes kann in vielerlei Hinsicht nützlich sein, ganz abgesehen davon, dass sich eine Detailvermessung am besten darauf stützen lässt.
Es wäre gewiss äußerst wichtig, wenn der größte Teil Europas vollständig von einem einzigen Netz überzogen wäre, und nach und nach werden wir das erreichen. Jeder Staat sollte es sich zur Ehre anrechnen, seinen Anteil daran so gut zu liefern, dass er würdig ist, neben den besten zu stehen.

\par\noindent\hspace*{\parindent}Ich glaube, Sie täten gut daran, auch den Meisner als Dreieckspunkt zu wählen. Sie werden ihn \textls[120]{unmittelbar} mit mindestens 5 meiner Dreieckspunkte verbinden können und wahrscheinlich zusätzlich mit ebenso vielen von \textsc{Müffling}.
Ich helfe Ihnen gern dabei und werde, sofern es sich irgendwie einrichten lässt, von meiner Seite aus dorthin beobachten.
Heliotrope erleichtern solche Arbeiten außerordentlich, und auf dem Meisner sind keine weiteren Vorkehrungen nötig, als dass ein steinernes Postament von etwa \(1\frac{3}{4}\) Fuß im Quadrat und 3 Fuß Höhe über dem Boden aufgestellt wird, um Heliotrop und Theodolit daraufzustellen, und dass vielleicht einige Bäume gefällt werden, falls sich sonst die
\par\clearpage
\noindent\null\hfill{\small\textbf{ZUR HANNOVERSCHEN TRIANGULATION.}}\hfill\makebox[0pt][r]{\fontsize{13}{16}\selectfont \textbf{381}}\par\medskip
\noindent freie Sicht nicht vollständig erreichen lässt.\quad\mbox{.\ .\ .\ .\ .}\quad Aus Hannover ist mir keinerlei Mitteilung über Ihre Triangulation zugegangen. Ich habe daher den Köterberg ausgeschlossen, der sonst vorzüglich zum genauen Anschluss des Hercules gedient hätte.
Inzwischen habe ich die Richtung der Linie Hohehagen-Hercules zu den anderen dortigen Linien sorgfältig gemessen.\\
\mbox{.\ .\ .\ .\ .\ .}
